\documentclass[12pt]{article}
\usepackage[utf8]{inputenc}
\usepackage[spanish]{babel}
\usepackage{amsmath}
\usepackage{amsfonts}
\usepackage{amssymb}
\usepackage{geometry}
\geometry{letterpaper, margin=2.5cm}

\begin{document}

\begin{center}
    \Large \textbf{Evaluación de Examen 1: Unidad I} \\
    \large Física para Ciencias de la Salud (005-1713) \\
    \vspace{0.5cm}
    \textbf{Estudiante:} María Fernanda Malavé \quad \textbf{C.I.:} 32.585.180
    \vspace{0.3cm} \\
    \large \textbf{Calificación Final: 10/10 puntos}
\end{center}

\vspace{0.5cm}

\section*{Análisis de Resultados}

\subsection*{1. Ángulo plano (Conversión a radianes)}
\textbf{Procedimiento y resultado:} Plantea correctamente la conversión usando una regla de tres. Muestra explícitamente la simplificación de la fracción dividiendo numerador y denominador entre 15, obteniendo de manera directa y metódica el resultado exacto de $\frac{5\pi}{12}$ rad. \\
\textbf{Veredicto:} \textbf{Correcto} (2/2 pts).

\subsection*{2. Sistemas de coordenadas (Longitud del hueso)}
\textbf{Procedimiento y resultado:} Excelente resolución analítica. La estudiante declara formalmente el cálculo de las componentes horizontales y verticales ($x_2-x_1$, $y_2-y_1$). Luego sustituye los valores en el Teorema de Pitágoras y desglosa las operaciones algebraicas ($\sqrt{36+64} = \sqrt{100}$) hasta obtener la longitud exacta de $10$ cm. \\
\textbf{Veredicto:} \textbf{Correcto} (2/2 pts).

\subsection*{3. Vectores (Fuerza resultante)}
\textbf{Procedimiento y resultado:} Realiza la suma de las componentes $\hat{i}$ y $\hat{j}$ de forma correcta y respetando las leyes de los signos, obteniendo las magnitudes correspondientes (30 y 60). Como detalle de formato, separa el vector con un punto y coma en lugar de usar un signo de suma ($30\hat{i} \text{ ; } 60\hat{j}$), sin embargo, los valores están perfectos. \\
\textbf{Veredicto:} \textbf{Correcto} (2/2 pts).

\subsection*{4. Potencias de diez (Notación científica y prefijo)}
\textbf{Procedimiento y resultado:} ¡Excelente resolución! Logró estructurar su respuesta cumpliendo ambos requerimientos: expresó la cifra en notación científica estricta ($1,25 \times 10^{-5}$ m) y adicionalmente ajustó el exponente a $10^{-6}$ para nombrar correctamente el prefijo del S.I. (``micro''). \\
\textbf{Veredicto:} \textbf{Correcto} (2/2 pts).

\subsection*{5. Sistemas de unidades (Conversión de densidad)}
\textbf{Procedimiento y resultado:} Aunque no utiliza una sola ecuación de factores en cadena, desglosa aritméticamente el problema en pasos lógicos impecables: divide entre 1000 para transformar la masa a kilogramos, y luego multiplica el resultado por $1.000.000$ para llevar el volumen a metros cúbicos. El resultado final numérico de $1030$ es exacto. \\
\textbf{Veredicto:} \textbf{Correcto} (2/2 pts).

\vspace{0.5cm}
\section*{Comentarios Generales y Sugerencias}
¡Felicidades por alcanzar la calificación máxima! El examen evidencia un excelente razonamiento analítico, orden procedimental y atención a los detalles (como se pudo notar al ser de las pocas estudiantes en contestar completamente la Pregunta 4). 
\begin{itemize}
    \item Como única sugerencia de formato, recordar que en álgebra vectorial, la notación estándar de componentes unitarios suele sumarse algebraicamente (Ej: $30\hat{i} + 60\hat{j}$) en lugar de separarse con punto y coma.
\end{itemize}

\end{document}